% SPDX-License-Identifier: Apache-2.0
% covers: Sd
\datasheet{Sd}{A native-mode SD card host on AXI-Lite: one command and
at most one block at a time, the block through a buffer of 128 words.}

\dsfacts{\code{lib/parts/src/sd.rs}}{\code{txhdl\_parts::sd::Sd}}%
{\code{//docs:examples}}

\subsection*{Function}

A card in SPI mode gives a few hundred kilobytes a second on one data
line, which is what \code{Spi} gives. The native mode is the card's
own protocol: a command line, \code{CMD}, that carries 48-bit commands
out and 48- or 136-bit responses back, and one or four data lines,
\code{DAT}, that carry blocks of 512 bytes either way, each line with
a CRC-16 of its own. Every command carries a CRC-7, and so does every
response except the one to \code{ACMD41}, whose CRC field is all
ones.

\code{Sd} is that host behind one AXI-Lite port, \code{bus}. A program
writes the argument and then the command word, and the command word
says what follows: no response, a short one or a long one; a block to
read or a block to write; whether to wait for the card's busy line
afterwards; and whether to check the response's CRC. The host sends
the command, takes the response, moves the block through its buffer,
which the program reads or fills a word at a time through
\code{data}, and says in \code{status} how it went.

A multi-block transfer is one command, \code{CMD18} or \code{CMD25},
then one data-only transfer for each further block, then
\code{CMD12}. The card's side is eight wires: \code{sclk}, the command
line in, out and its output enable, the four data lines in, out and
their output enable, and \code{irq}.

\subsection*{Parameters}

None.

\subsection*{Register map}

The word is address bits 5 to 2. A word the map does not name reads as
zero.

\dsregs{Sd}

\code{cmd}'s \code{resp} is 0 for none, 1 for short and 2 for long.
Writing \code{cmd} starts the command. Writing a one to
\code{status}'s \code{done} clears it and the four faults with it.

A half of a card clock takes \code{div + 1} cycles, so the card's
clock is the system's over \code{2 * (div + 1)}: at 100~MHz, 124 is
the 400~kHz a card starts at, 1 is 25~MHz and 0 is 50~MHz.

\dstables{Sd}

\subsection*{Behaviour}

\begin{itemize}
\item \textbf{A write to \code{cmd} while a command runs is ignored.}
A program reads \code{status} first.
\item The host drives its lines on the falling edge of the clock it
makes and samples the card's on the rising edge. The clock runs only
while a command runs, and is low otherwise.
\item A response that has not started 64 card clocks after the command
ends is a timeout. A block, the CRC status token after a written
block, and the end of the card's busy are each given $2^{24} - 1$ card
clocks.
\item A CRC is checked the way a shift register checks one: the
received CRC goes into the same register after the data, and what is
left is zero when the two agree. For a long response the CRC covers
the 120 bits after the header. With four lines each line has its own
CRC-16, and any of the four being wrong sets bit 5.
\item A written block is accepted only when the card's CRC status
token is \code{010}; anything else sets bit 5. A command whose
response CRC is wrong goes no further: its block does not move.
\item Starting a read sets the write pointer to zero, and starting a
write sets the read pointer to zero, so a block fills or empties the
buffer from its start. A read of \code{data} moves the read pointer
and a write moves the write pointer, each modulo 128.
\item \code{irq} is done and the interrupt enable together, so it
stays high until the program clears done.
\item Every AXI-Lite access is answered \code{OKAY}.
\end{itemize}

\subsection*{Verification}

\begin{itemize}
\item \code{ex\_sd} does what a driver does, through an AXI4 host and
\code{LiteBridge<1, ..>}, against \code{SdCard}, the model of a card in \code{sd.rs},
stepped on the wires every cycle. It brings the card up with
\code{CMD0}, \code{CMD8}, \code{ACMD41}, \code{CMD2}, \code{CMD3},
\code{CMD7} and \code{ACMD6}, reads block 3 on one line, then writes
block 5 on four and reads it back, and checks both blocks.
\item The netlist is checked against that run under nvc and
Verilator: \code{//docs:sd\_sim\_sd\_tb\_test} and
\code{//docs:sd\_vsim\_test}.
\item \code{//lib/parts:parts\_test} holds six tests of it:
\begin{itemize}
\item the card comes up and gives its CID;
\item a block is read on one line;
\item a block is written and read back on four lines, the card's CRC
status accepting it;
\item two blocks each way, as a command, a data-only transfer and
\code{CMD12};
\item a spoiled response CRC sets bit 3 and the next command is
clean; a command the card does not answer sets bit 2; a spoiled block
sets bit 5;
\item the CRC-7 of \code{CMD0} with a zero argument is \code{0x4a} and
the CRC-16 of 512 bytes of \code{0xff} is \code{0x7fa1}, the
specification's own examples.
\end{itemize}
\end{itemize}

\subsection*{Used by}

The example \code{ex\_sd}. Nothing on a board holds it yet.

\subsection*{Limits}

\textbf{Not yet on hardware.} The board half of issue 153, the card
slot, reading the card's CID and block 0 from a real card, is owed,
and nothing here has met a real card.

One command and at most one block at a time, and the block moves
through a register. A card streams a multi-block read back to back,
and keeping up with that at 25~MHz a word at a time is what direct
memory access would be for; there is none.

\code{SdCard}, the model, is not a component: nothing lowers it, and
it is only as right as the tests that check the host against it.

It runs on one clock, \code{DefaultClock}.
